% ---------------------------------------------------------------- Limitations
\section*{Limitations}
\begin{itemize}
\item \textbf{Scale of the workloads.} All runs are small image classifiers trained on 4{,}000-image subsets for 40 epochs on a CPU. Large models, long schedules with warm-up and decay, and fine-tuning of pretrained networks produce curve shapes that are not represented here; a gate for such jobs should be retrained on their curves (Steps 1 to 4), and its transfer should be checked with the leave-one-workload-out protocol before it is trusted.
\item \textbf{No accuracy advantage over simpler rules.} On the same hindsight labels a linear gate reaches a similar trade-off, a patience rule tuned on other workloads is as good on the stop decision alone, and in live training a reduce-on-plateau schedule with patience-based stopping was at least as good as the full controller (Method validation). The contribution of the KAN is an exact, readable and non-linear decision rule, together with the training signal that makes such a rule useful; it is not higher accuracy. Where a linear rule is sufficient for a workload, it is also acceptable and even simpler to audit.
\item \textbf{The learning-rate action is not learned.} The adjust band is derived from the stop score, not from the measured effect of a learning-rate change. In live control it did not outperform a reduce-on-plateau schedule and harmed some slowly improving runs, so the gate is best used as the stop rule together with a standard learning-rate schedule (Method validation). Learning when a learning-rate change helps would need labels from paired runs with and without the change.
\item \textbf{Offline replay covers stopping only.} A recorded curve cannot show what a learning-rate change would have done, so the offline benchmark scores the stop decision and treats \emph{adjust} as \emph{continue}. Learning-rate adjustment is evaluated only in the closed-loop experiments, on 12 runs per condition.
\item \textbf{Hindsight labels need complete runs.} Runs that a gate itself stopped have no observed future and cannot be labelled; harvesting labels from a deployed gate is therefore right-censored. Collecting a share of uncontrolled runs, or runs continued past the gate's stop for a few epochs, avoids this bias.
\item \textbf{Budget feature and budget length.} The budget feature is the fraction of the epoch budget remaining, not measured accelerator time, and all gates were trained and tested with a 40-epoch budget. Because $x_5$ is normalised by a fixed 30 epochs, a gate should be retrained, or at least re-validated with the leave-one-workload-out protocol, before it is used with a much longer budget or on heterogeneous hardware.
\item \textbf{Validation and test data.} Stopping decisions and reported quality both use the validation split, which here is drawn from the official test split. In an application with a separate test set, the gate must read a validation split carved from the training data, and the test set is used only once, after training has stopped.
\item \textbf{Noisy validation signals.} The decision bands are 35 points wide, but a score close to a band boundary can flip between decision points. With very noisy validation losses, requiring the stop band on two consecutive decision points is a simple hysteresis; it was not needed or evaluated here.
\item \textbf{Supervisor agent not evaluated.} All experiments ran with the rule-based fallback supervisor (no LLM). Because the gate is binding, the supervisor's proposal does not change any action; its rationales were not evaluated.
\item \textbf{Cluster deployment.} The released controller includes a Kubeflow Trainer client, but all reported runs used the local process client. Behaviour on a live Kubernetes cluster, such as scheduling delays between a decision and its effect, was not measured.
\end{itemize}
